% Assembly lemma for the actual bad-semiprime count, 2026-10-08.
% Does not certify its two analytic inputs or edit the main manuscript.
\subsection{From a compiled upper weight to the actual bad count}
\label{sec:actual-k-positive-assembly}

Put $\mathcal H_N=C_NN/\log^2N$, with the same singular series as
the manuscript, and let $N$ tend to infinity through even integers.
The following argument isolates exactly what is needed on the modulus
side. In particular, a rational density bound alone is not its hypothesis.

\paragraph{The two inputs.}
Fix $K_0>0$, $0<\theta<1$, and $0<\rho<1$. Suppose a fixed finite
compiled strategy supplies real coefficients $\rho_N(d)$ such that
\[
 \operatorname{supp}\rho_N\subset
 \{d\le N^\theta:d\mid P_N(N^\rho)\},\qquad |\rho_N(d)|\le1,
\]
and, for every integer $n$,
\begin{equation}\label{eq:actual-k-majorant-input}
 \sum_{d\mid n}\rho_N(d)\ge
 \mathbf1_{(n,P_N(N^\rho))=1},\qquad
 \sum_d\frac{\rho_N(d)}{\varphi(d)}
 \le (K_0+o(1))\frac{C_N}{\log N}.
\end{equation}
All supported $d$ are squarefree and coprime to $N$. The first
inequality includes $n=0$; its left side is a finite sum.

The second input is the actual signed distribution statement for
each nonnegative smooth two-prime weight $W_N(r,q)$ used below,
with a fixed positive factor cutoff $r,q\ge N^a/2$:
\begin{equation}\label{eq:actual-k-distribution-input}
 \sum_d\rho_N(d)\left\{
 \sum_{rq\equiv N\ (d)}W_N(r,q)
       -\frac{X_W^{(d)}}{\varphi(d)}\right\}
 =O_{A,W}(N\log^{-A}N),
\end{equation}
where sums over $r,q$ are over primes,
$X_W^{(d)}=\sum_{(rq,d)=1}W_N(r,q)$, and $X_W=\sum W_N(r,q)$.
Weights have bounded size, product $rq\ll N$, and bounded normalized
smooth derivatives for each fixed smoothing precision. The constants
may depend on this precision and the fixed strategy but not on $N$.
These inputs must hold for the \emph{same} coefficients $\rho_N$.
They are the tasks of the compilation/density and distribution audits,
respectively, and are not deduced here from their saved status strings.

\paragraph{Exact coprime mass can be removed after summing the moduli.}
For a prime $p$, squarefreeness and the support bound give
\[
 \sum_{d:p\mid d}\frac{|\rho_N(d)|}{\varphi(d)}
 \le\frac1{p-1}\sum_{e\le N^\theta/p}\frac1{\varphi(e)}
 \ll\frac{\log N}{p-1}.
\]
Here $\sum_{e\le T}1/\varphi(e)\ll\log(2T)$ is elementary: use
$n/\varphi(n)=\sum_{h\mid n}\mu^2(h)/\varphi(h)$ and the convergent
Euler product for $\sum_h\mu^2(h)/(h\varphi(h))$.
Positivity of $W$ and a union bound over its two prime slots therefore
yield
\begin{align*}
 \left|\sum_d\frac{\rho_N(d)}{\varphi(d)}
                     (X_W^{(d)}-X_W)\right|
 &\le\sum_{r,q}W_N(r,q)
       \sum_{p\in\{r,q\}}\sum_{d:p\mid d}
                                \frac{|\rho_N(d)|}{\varphi(d)}\\
 &\ll_a X_W N^{-a}\log N=o(\mathcal H_N).
\end{align*}
The bound remains valid on the diagonal $r=q$; double-counting there
only weakens an upper bound. There is no factor $N^\theta$ in this
error. A fixed positive lower bound for $C_N$ is sufficient throughout.

\paragraph{Positive majorization is used before distribution.}
If $N-rq$ is a prime at least $N^\rho$, its upper-sieve indicator is
one. The number of target pairs with $0<N-rq<N^\rho$ is
$O(N^\rho)$: for each integer $rq=N-p$, there is at most one ordered
prime factorization $r\le q$. Thus for a majorant of an ordered domain,
\[
 \#\{(r,q)\text{ in the domain}:N-rq\text{ prime}\}
 \le O(N^\rho)+\sum_{r,q}W_N(r,q)\sum_{d\mid N-rq}\rho_N(d).
\]
An outward smoothing may include $rq>N$ or $rq=N$. This causes no
sign problem: the pointwise divisor sum is nonnegative at every
integer, including zero. No estimate on a signed boundary discrepancy
is being inferred from its small mass.

The radial majorant can also include products below the local scales
on which the distribution proof is applied. These are removed with a
power saving, rather than by postulating distribution at those scales.
For a fixed $\kappa>0$, their absolute progression contribution is
\[
 \sum_{m<2N^{1-\kappa}} O_W(\tau_C(m))
              \sum_{d\mid N-m}|\rho_N(d)|
 \ll N^{1-\kappa+o(1)},
\]
since $m\ne N$ and fixed divisor orders are $N^{o(1)}$ pointwise.
Their exact-coprime mean costs at most the same mass times
$\sum_{d\le N^\theta}1/\varphi(d)\ll\log N$.
One may insert a smooth product cutoff equal to zero for
$m\le N^{1-\kappa}$ and one for $m\ge2N^{1-\kappa}$.
Its multiplicative derivatives are uniformly bounded on dyadic
charts; the entire transition is covered by the preceding estimate.
Thus the global radial majorant needs distribution only on the
retained near-$N$ product scales, with no factor $N^\theta$ loss.

Apply \eqref{eq:actual-k-distribution-input}, remove the exact coprime
mass as above, and use \eqref{eq:actual-k-majorant-input}. If
$X_W=(J_W+o(1))N/\log N$, the resulting bound is
\[
 (K_0J_W+o(1))\mathcal H_N.
\]

\paragraph{Mass and smoothing of the bad-semiprime domain.}
Fix $1/2<b<1$, put $\beta=b/(1+b)$, and choose sufficiently small
fixed $0<\kappa<1/2$ as required by the distribution input.
Products $rq<N^{1-\kappa}$ contribute at most $N^{1-\kappa}$ to
the actual ordered count. On the remaining bad domain
\[
 rq<N,\qquad r\le q,\qquad r>q^b,
\]
both factors exceed $N^{(1-\kappa)\beta}$. In particular the fixed
cutoff $a=\beta/4$ is available with room for smoothing.

Use $t=\log r/\log N$ and $y=rq/N$. The limiting labelled
prime-product measure is
\[
 \frac{N}{\log N}\frac{dt\,dy}{t(1-t)}.
\]
For $y\ge\eta>0$ this follows by first applying PNT in
$q=Ny/r$ and then partial summation to the reciprocal-prime sum in
$r$. The accumulated PNT error is $O_A(N\log^{-A}N)$ because
$\sum_{r\ge N^a/2}1/r=O_a(1)$ over the relevant primes.
For $y<\eta$, the elementary prime upper bound gives total mass
$O_a(\eta N/\log N)$ by the same reciprocal-prime sum.
Hence $N\to\infty$ followed by $\eta\downarrow0$ is valid.

The exact exponent inequalities are
\[
 t\le\frac12+\frac{\log y}{2\log N},\qquad
 t>\beta+\frac{b\log y}{(1+b)\log N}.
\]
Thus the limiting domain is $\beta<t\le1/2$, $0<y<1$, and
\[
 J_b=\int_0^1\!dy\int_\beta^{1/2}\frac{dt}{t(1-t)}
     =\log(1/b).
\]
The boundary $rq=N$ is smoothed in $y$, not declared null in the
full exponent plane. The slanted exponent walls are null on the
one-dimensional simplex at each $y>0$. The positive approximation
of \texttt{prime-product-radial-boundary-repair.tex} therefore gives,
for every fixed $\varepsilon>0$, a smooth majorant with
$J_W\le J_b+\varepsilon$. The fixed lower radial cutoff
$N^{-\kappa}$ has vanishing limiting mass. Repeated factors give
$O(\sqrt N)$ possible pairs and do not affect $J_b$.

Apply the positive bound for this majorant, let $N\to\infty$, and
then let $\varepsilon\downarrow0$. We conclude from the two stated
inputs that
\begin{equation}\label{eq:actual-k-bad-count-conclusion}
 \limsup_{\substack{N\to\infty\\2\mid N}}
 \frac{B_b(N)}{\mathcal H_N}\le K_0\log(1/b).
\end{equation}
All auxiliary choices are fixed before the $N$ limit; no averaging
over even $N$ and no effective threshold is asserted.

\paragraph{Endpoint slack and the final logical status.}
Suppose the joint tree audit proves that every $K_0>K_{\rm end}$
admits a common finite buffered strategy satisfying both inputs.
Take any $K_{\rm end}<K_0<67619/10000$. Then
\eqref{eq:actual-k-bad-count-conclusion} proves precisely the
upper-count statement required by
\texttt{external-published-chen-lower-k-independent-transfer.tex}.
It is unnecessary to attain $K_{\rm end}$ with a single finite tree.
The rational slack in that fragment subsequently absorbs all $o(1)$
terms in its $b=747/1000$ deduction. This assembly proof establishes
the implication, not the outstanding hypotheses or the validity of
any external Chen input.
